\chapter{Supersymmetry as Odd Symmetry and Superspace}

In Leonard Susskind's fifth lecture on supersymmetry and grand unification, curated by LazyingArt LLC, the subject begins with a warning. Supersymmetry is not introduced as a pleasant variant of familiar symmetry, but as a remarkably abstract extension of space-time symmetry, and one does not even state it cleanly without first learning how to manipulate Grassmann variables. These notes therefore follow the lecture's own progression. We begin with the odd and even algebra, move through the small calculus of Grassmann variables, pause to recall the ordinary algebra of continuous symmetries, and only then construct the simplest supersymmetry algebra. The real surprise is that the anticommutator of the odd generators does not close on another internal generator; it closes on the Hamiltonian. From there the lecture turns geometric and points toward superspace.

\section{Why Supersymmetry Needs a Strange Prelude}

Supersymmetry is presented here as a symmetry of space-time, but of a very unfamiliar kind. Rotations, translations, and the ordinary continuous symmetries of physics act on ordinary coordinates or on internal labels. Supersymmetry does something stranger: it mixes bosonic and fermionic degrees of freedom, and in doing so it forces us to enlarge the very notion of space-time.

That is why the lecture does not begin with a direct definition of the supersymmetry algebra. It first backs up and reviews Grassmann numbers. This is not a side topic. The odd generators of supersymmetry are built from the same sort of odd objects, and if one is not comfortable with nilpotent variables, anticommuting signs, and the calculus attached to them, the later formulas look like pure magic. Susskind's point is that we should first learn the bookkeeping before we ask the bookkeeping to carry physical meaning.

\section{Grassmann Numbers, Functions, and Calculus}

\paragraph{Odd and even elements.}
Let \(\epsilon_i\) denote Grassmann variables. They are odd elements, and they anticommute:
\begin{equation}
\epsilon_i \epsilon_j + \epsilon_j \epsilon_i = 0.
\end{equation}
In particular,
\begin{equation}
\epsilon_i^2 = 0.
\end{equation}
Ordinary numbers, real or complex, are even elements. The parity rules are
\begin{equation}
\text{odd}\cdot\text{odd}=\text{even},
\qquad
\text{odd}\cdot\text{even}=\text{odd},
\qquad
\text{even}\cdot\text{even}=\text{even}.
\end{equation}
Even elements commute with everything. Odd elements anticommute only with other odd elements.

\paragraph{Functions terminate.}
A function of one Grassmann variable is only a short polynomial,
\begin{equation}
f(\epsilon)=a+b\,\epsilon,
\end{equation}
because \(\epsilon^2=0\). With two independent odd variables,
\begin{equation}
f(\epsilon_1,\epsilon_2)=a+b_1\epsilon_1+b_2\epsilon_2+c\,\epsilon_1\epsilon_2.
\end{equation}
Nothing higher survives. This is one of the reasons the lecture keeps returning to the word bookkeeping: there is structure here, but it is a small and highly constrained structure.

If we organize two real odd variables into a complex pair \(\theta,\bar\theta\), then the same finite expansion can be written schematically as
\begin{equation}
f(\theta,\bar\theta)=a+\bar b\,\theta+\bar\theta\,b+c\,\bar\theta\theta.
\end{equation}
This is the natural form suggested by the lecture, although the detailed sign bookkeeping depends on the parity assignments of the coefficients and on the chosen ordering convention.

\paragraph{Differentiation.}
The derivative with respect to a Grassmann variable is simple once one respects ordering. If a function does not depend on \(\theta\), then
\begin{equation}
\partial_\theta a = 0.
\end{equation}
If the dependence is written with \(\theta\) adjacent to the derivative in the standard way, then
\begin{equation}
\partial_\theta(a+\theta b)=b.
\end{equation}
But if an odd coefficient sits to the left, the sign changes when we reorder. For example, if \(b\) is odd, then \(b\theta=-\theta b\), so
\begin{equation}
\partial_\theta(b\theta)=-b.
\end{equation}
This is the rule Susskind insists on: before differentiating, bring the differentiated Grassmann variable next to the derivative, and account for every sign picked up while moving odd elements past odd elements.

\paragraph{A worked identity.}
The lecture then makes an important comparison with ordinary calculus. For an ordinary variable \(x\),
\begin{equation}
[\partial_x,x]=1.
\end{equation}
For a Grassmann variable, the corresponding identity uses an anticommutator. Let
\begin{equation}
f(\theta)=a+\theta b.
\end{equation}
Then
\begin{align}
\theta\,\partial_\theta f &= \theta b, \\
\partial_\theta(\theta f) &= \partial_\theta(\theta a)=a,
\end{align}
because \(\theta^2=0\). Adding the two actions gives
\begin{equation}
\bigl(\partial_\theta\theta+\theta\partial_\theta\bigr)f
=
a+\theta b
=
f.
\end{equation}
Since this works on an arbitrary test function of one Grassmann variable,
\begin{equation}
\{\partial_\theta,\theta\}=1.
\end{equation}
This is the first place where the lecture turns Grassmann calculus into operator algebra. We are no longer merely manipulating nilpotent symbols; we are already seeing the kind of relations from which symmetry generators are built.

\paragraph{Integration.}
Susskind only states the integral rule, but it is worth recording because he uses it to emphasize how small the Grassmann calculus really is:
\begin{equation}
\int d\theta\, f(\theta)=\partial_\theta f(\theta).
\end{equation}
He motivates this by remarking that if one wants the integration-by-parts rule to survive, then derivative and integral must coincide for Grassmann variables.

\subsection{Question \& Answer}

A natural question arises as soon as we write a term such as \(c\,\bar\theta\theta\): does the order of \(\theta\) and \(\bar\theta\) matter in any deep way?

At this stage the answer is no, not in the sense of changing the parity of the object. The order matters only through signs. Since \(\theta\) and \(\bar\theta\) are both odd,
\begin{equation}
\bar\theta\theta=-\theta\bar\theta.
\end{equation}
Thus one may equally well write
\begin{equation}
c\,\bar\theta\theta = -\,c\,\theta\bar\theta,
\end{equation}
or absorb the sign into the coefficient. If \(c\) is itself odd, further sign changes appear when \(c\) is moved past odd variables. But none of these manipulations changes odd into even or even into odd. The lecture's answer is deliberately modest: at this point the signs are a bookkeeping device for tracking fermionic structure.

\section{Ordinary Symmetry Algebras as the Comparison Case}

Before defining supersymmetry, the lecture pauses to recall what an ordinary continuous symmetry looks like. Let \(G_i\) be the infinitesimal generators of a continuous group. Their algebra closes under commutators,
\begin{equation}
[G_i,G_j]=i\,c_{ijk}G_k,
\end{equation}
with \(c_{ijk}\) the structure constants. This closure property is the infinitesimal version of group multiplication.

The generators of sensible symmetry groups are Hermitian operators. One may nevertheless combine them into useful complex linear combinations. Angular momentum provides the familiar example:
\begin{equation}
L_\pm = L_x \pm i L_y.
\end{equation}
The underlying observables remain Hermitian even when one passes to raising and lowering combinations.

For a genuine symmetry of a physical system, there is one more condition:
\begin{equation}
[H,G_i]=0.
\end{equation}
This means that the action of the generator does not change the energy. Equivalently, the generator is conserved. The lecture wants this conventional picture fresh in mind, because supersymmetry will look almost the same at first glance and profoundly different a moment later.

One further contrast matters. The ordinary observables of physics are even objects. Measured quantities return ordinary numbers, not odd Grassmann elements. Supersymmetry will break precisely that habit: its generators are odd.

\section{A Minimal Supersymmetry Algebra from Oscillators}

To make the new algebra concrete, Susskind constructs the most primitive model he can: one bosonic oscillator and one fermionic oscillator, both thought of at rest. The bosonic creation and annihilation operators \(a_+\) and \(a_-\) satisfy
\begin{equation}
[a_+,a_+]=[a_-,a_-]=0,
\qquad
[a_-,a_+]=1.
\end{equation}
The fermionic creation and annihilation operators \(c_+\) and \(c_-\) satisfy
\begin{equation}
\{c_+,c_+\}=\{c_-,c_-\}=0,
\qquad
\{c_-,c_+\}=1.
\end{equation}
In particular,
\begin{equation}
c_+^2=c_-^2=0.
\end{equation}
The bosonic operators are even, the fermionic operators are odd, and every \(a\) commutes with every \(c\).

The odd generators of the toy supersymmetry are then defined by
\begin{equation}
Q=\sqrt{m}\,a_+c_-,
\qquad
\bar Q = Q^\dagger = \sqrt{m}\,a_-c_+.
\end{equation}
The lecture uses \(\bar Q\) and \(Q^\dagger\) interchangeably. The action is simple: \(Q\) removes a fermion and replaces it by a boson, while \(\bar Q\) does the reverse.

Now comes the key calculation. We compute
\begin{align}
\{Q,\bar Q\}
&=
m\bigl(a_+c_-a_-c_+ + a_-c_+a_+c_-\bigr) \\
&=
m\bigl(a_+a_-c_-c_+ + a_-a_+c_+c_-\bigr),
\end{align}
where we have commuted the even \(a\)'s past the odd \(c\)'s. Next use
\begin{equation}
c_-c_+ = 1-c_+c_-,
\qquad
a_-a_+ = a_+a_-+1.
\end{equation}
Then
\begin{align}
\{Q,\bar Q\}
&=
m\Bigl[a_+a_-\bigl(1-c_+c_-\bigr) + \bigl(a_+a_-+1\bigr)c_+c_-\Bigr] \\
&=
m\bigl(a_+a_- + c_+c_-\bigr).
\end{align}
If we define the boson and fermion number operators by
\begin{equation}
N_B=a_+a_-,
\qquad
N_F=c_+c_-,
\end{equation}
then
\begin{equation}
\{Q,\bar Q\}=m(N_B+N_F).
\end{equation}
But in this simple model the total rest energy is just the common mass times the total number of particles, so
\begin{equation}
H=m(N_B+N_F).
\end{equation}
Therefore,
\begin{equation}
\{Q,\bar Q\}=H.
\end{equation}

\begin{figure}
\centering
\includegraphics[width=0.62\textwidth]{lecture_05_figure_02.png}
\caption{Supersymmetry anticommutator yielding the Hamiltonian.}
\end{figure}

This is the first real climax of the lecture. The anticommutator of two odd generators does not simply return another internal generator. It returns the Hamiltonian, and therefore something tied directly to time translation.

The rest of the algebra is immediate:
\begin{equation}
\{Q,Q\}=0,
\qquad
\{\bar Q,\bar Q\}=0,
\end{equation}
because \(Q^2\) contains \(c_-^2\) and \(\bar Q^2\) contains \(c_+^2\).

\subsection{Question \& Answer}

Why does \(Q\) commute with \(H\), and why must the boson and fermion masses match?

The lecture answers this by acting on an energy eigenstate. Let \(|E\rangle\) be any state of definite energy:
\begin{equation}
H|E\rangle = E|E\rangle.
\end{equation}
If the boson and fermion masses are equal, then \(Q\) replaces one particle by another of exactly the same mass. Therefore \(Q|E\rangle\) is again a state with the same energy \(E\). On such a state,
\begin{align}
(QH-HQ)|E\rangle
&=
Q(E|E\rangle)-H(Q|E\rangle) \\
&=
E\,Q|E\rangle-E\,Q|E\rangle \\
&=0.
\end{align}
If the energy eigenstates form a complete basis, this implies
\begin{equation}
[Q,H]=0.
\end{equation}
The same reasoning applies to \(\bar Q\). Thus the odd generators are conserved, exactly as symmetry generators should be.

But the equal-mass assumption is essential. If the boson and fermion masses differed, the state \(Q|E\rangle\) would not have the same energy as \(|E\rangle\), the commutator would no longer vanish, and the neat closed algebra would be spoiled. This is the point in the lecture where the algebraic claim is slowed down and justified at the level of states.

\section{From Time Translation to Superspace}

Once the Hamiltonian has appeared on the right-hand side of the supersymmetry algebra, the lecture backs up and asks what the Hamiltonian is doing there. The answer is that \(H\) is the generator of time translations.

For a state vector, the infinitesimal time-evolution law is the usual Schr\"odinger equation,
\begin{equation}
i\,\frac{\partial}{\partial t}|\psi\rangle = H|\psi\rangle.
\end{equation}

\begin{figure}
\centering
\includegraphics[width=0.56\textwidth]{lecture_05_figure_03.png}
\caption{Time derivative acting on a state vector.}
\end{figure}

A finite shift of time by an amount \(\delta\) may be written as
\begin{equation}
e^{-iH\delta}|\psi(t)\rangle = |\psi(t+\delta)\rangle.
\end{equation}
The lecture then translates the same idea into the language of a function. If we promote the object under discussion to
\begin{equation}
\psi(\theta,\bar\theta,t),
\end{equation}
then an ordinary time shift does nothing to \(\theta\) and \(\bar\theta\), but it changes the function by
\begin{equation}
t\to t+\delta,
\qquad
\psi(\theta,\bar\theta,t)\to
\psi(\theta,\bar\theta,t)+\delta\,\frac{\partial \psi(\theta,\bar\theta,t)}{\partial t}.
\end{equation}

\begin{figure}
\centering
\includegraphics[width=0.72\textwidth]{lecture_05_figure_04.png}
\caption{First-order time translation of the wavefunction.}
\end{figure}

So far this is ordinary time-translation invariance. The exotic step is to ask whether one can invent a meaningful symmetry that shifts the Grassmann coordinate itself. Susskind proposes the one-sided transformation
\begin{equation}
\theta \to \theta+\epsilon,
\qquad
t \to t+i\,\bar\theta\,\epsilon,
\end{equation}
with \(\epsilon\) an odd parameter. Because \(\epsilon^2=0\), such a shift is automatically infinitesimal in the Grassmann sense.

The induced change in \(\psi\) is first order in \(\epsilon\). A conventional reconstruction of the operator suggested by the lecture is
\begin{equation}
\delta_\epsilon \psi
=
\epsilon\Bigl(\partial_\theta - i\,\bar\theta\,\partial_t\Bigr)\psi.
\end{equation}
This motivates the odd differential operator
\begin{equation}
q = \partial_\theta - i\,\bar\theta\,\partial_t.
\end{equation}
By complex conjugation one introduces a partner
\begin{equation}
\bar q = \partial_{\bar\theta} + i\,\theta\,\partial_t.
\end{equation}
Here one should be cautious. Susskind explicitly says he may have misplaced a sign and probably a factor of two, so these formulas should be read as a conventional normalization of the construction rather than as a literal transcription of every board sign.

What matters is the algebraic upshot. The mixed terms in the anticommutator survive because of the Grassmann identities
\begin{equation}
\{\partial_\theta,\theta\}=1,
\qquad
\{\partial_{\bar\theta},\bar\theta\}=1,
\end{equation}
while the derivative-derivative and odd-coordinate--odd-coordinate terms vanish. Consequently,
\begin{equation}
\{q,\bar q\}\propto -\,i\,\partial_t.
\end{equation}
But \(-i\partial_t\) is the generator of time translation, hence in this simple model it is the Hamiltonian. This is the geometric rewriting of the earlier oscillator algebra. Supersymmetry generators become translations in odd directions of an enlarged space, and their anticommutator produces an ordinary translation in time.

This enlarged arena is what the lecture calls superspace. Ordinary motion moves us in \(x\) and \(t\); supersymmetry adds motion in \(\theta\) and \(\bar\theta\) as well. That is the conceptual endpoint of the lecture.

\section{Summary}

The lecture builds supersymmetry in two passes. First it teaches us how to live with odd variables: anticommuting numbers, truncated functions, derivatives that care about order, and the operator identity \(\{\partial_\theta,\theta\}=1\). Then it returns to the familiar language of symmetry generators and changes one decisive rule: odd generators close by anticommutators.

In the minimal oscillator model, that rule produces
\begin{equation}
\{Q,\bar Q\}=H,
\end{equation}
so the algebra closes only when the Hamiltonian is included. This is what makes the symmetry new. It is not merely an internal transformation taking fermions to bosons; done twice in the appropriate anticommuting sense, it generates time translation. In one sense that is still bookkeeping for fermions and bosons. In another sense it is a much deeper enlargement of the geometry on which the theory lives.